\documentclass[11pt]{article}
\usepackage[utf8]{inputenc}
\usepackage[T1]{fontenc}
\usepackage[italian]{babel}
\usepackage{geometry}
\geometry{a4paper, margin=1in}
\usepackage{xcolor}
\usepackage{listings}
\usepackage{hyperref}

\usepackage{tcolorbox}
\tcbuselibrary{breakable}
\begin{document}

\begin{tcolorbox}[colback=blue!5!white,colframe=blue!75!black,title=\textbf{DOMANDA}, breakable]
Temi d'esame (per ognuno ragionare sulla scelta del protocollo di livello trasporto)
\begin{itemize}
    \item Implementare una semplice versione di terminale remoto. Il client si connette ad un server di cui si conosce IP e porta e chiede all'utente di scrivere un comando come in un normale terminale a caratteri (detto anche shell o console) di Linux. Il comando viene eseguito dal server e il relativo output viene stampato a video dal client.
\end{itemize}
\end{tcolorbox}

\begin{tcolorbox}[colback=green!5!white,colframe=green!75!black,title=\textbf{RISPOSTA}, breakable]
\textbf{Piano di Implementazione:} Terminale Remoto (es20)

\textbf{1. Scelta del Protocollo di Trasporto}\\
La traccia richiede che il client si connetta a un server, invii un comando (es. comando di shell) e riceva l'output risultante dall'esecuzione di quel comando dal server.

\textbf{Scelta raccomandata:} \texttt{TCP}

\textbf{Motivazione:}\\
In questo scenario è fondamentale che i dati non vadano persi né vengano ricevuti in ordine errato. Un comando potrebbe risultare corrotto (ad es. un "ls -l" potrebbe arrivare incompleto), ma l'aspetto critico è la ricezione dell'output. L'output di un comando di shell può essere di dimensioni arbitrarie e superare ampiamente il limite di singolo pacchetto (\texttt{MTU}) supportato tipicamente da \texttt{UDP}. Il protocollo \texttt{TCP} garantisce una trasmissione orientata al flusso (stream) e gestisce nativamente la segmentazione e la ritrasmissione dei pacchetti qualora essi si perdano nella rete.

Il server leggerà l'output tramite una pipe (ad es. usando \texttt{popen} in C) e lo invierà sul socket \texttt{TCP} in vari blocchi. Il client continuerà semplicemente a leggere dal socket finché non riceverà un segnale di termine (la chiusura della connessione o fine dei dati). Con \texttt{UDP}, al contrario, sarebbe necessario implementare a livello applicativo la frammentazione dell'output e gestire esplicitamente ACK/NACK per garantire che il client riceva tutto l'output ordinato correttamente.

\textbf{2. Implementazione}
\begin{itemize}
    \item Il client \texttt{TCP} legge il comando inserito dall'utente (es. tramite \texttt{fgets}), lo invia al server tramite \texttt{TCPSend}, e poi entra in un ciclo in cui legge (\texttt{TCPReceive}) l'output inviato dal server e lo stampa a video, fermandosi quando la trasmissione è completata (chiusura del socket o invio terminato).
    \item Il server \texttt{TCP} accetta connessioni. Quando riceve un comando, usa la chiamata di sistema "\texttt{popen(comando, "r")}" per eseguirlo come comando di shell e aprire una pipe per leggerne l'output. Legge l'output a blocchi usando \texttt{fread} e li invia man mano al client tramite \texttt{TCPSend}. Finito l'output, chiude il file associato a \texttt{popen} (\texttt{pclose}) e chiude il socket, il che indicherà al client che il risultato del comando è stato interamente trasmesso.
\end{itemize}
\end{tcolorbox}

\end{document}
